\begin{frame}{Bubble Sort의 Mission과 Invariant}
\mission{인접한 역전 쌍을 교환하여 큰 원소를 오른쪽으로 떠오르게 한다.}
\begin{block}{Invariant}왼쪽부터 한 pass를 마치면 active 구간의 최댓값이 오른쪽 끝에 있다.\end{block}
\end{frame}
\begin{frame}{조기 종료 Bubble Sort}
\begin{algorithmic}[1]\Procedure{BubbleSort}{$A[1..n]$}\For{$last\gets n$ downto $2$}\State $swapped\gets\mathrm{false}$\For{$i\gets1$ to $last-1$}\If{$A[i]>A[i+1]$}\State \Call{Swap}{$A[i],A[i+1]$}; $swapped\gets\mathrm{true}$\EndIf\EndFor\If{not $swapped$}\State \Return\EndIf\EndFor\EndProcedure\end{algorithmic}
\end{frame}
\begin{frame}{Bubble Sort: 인접 비교의 흐름}
\optional[추가 예제]\hfill
\centering\begin{tikzpicture}[x=1.18cm]
\only<1>{\def\vals{{29,10,14,37,13}}\def\a{1}\def\b{2}\def\swap{1}}
\only<2>{\def\vals{{10,29,14,37,13}}\def\a{2}\def\b{3}\def\swap{1}}
\only<3>{\def\vals{{10,14,29,37,13}}\def\a{3}\def\b{4}\def\swap{0}}
\only<4>{\def\vals{{10,14,29,37,13}}\def\a{4}\def\b{5}\def\swap{1}}
\only<5>{\def\vals{{10,14,29,13,37}}\def\a{0}\def\b{0}\def\swap{0}}
\foreach \i in {1,...,5}{\pgfmathparse{\vals[\i-1]}\edef\v{\pgfmathresult}\ifnum\i=\a\ifnum\swap=1\node[sort current]at(\i,0){\v};\else\node[sort compared]at(\i,0){\v};\fi\else\ifnum\i=\b\ifnum\swap=1\node[sort current]at(\i,0){\v};\else\node[sort compared]at(\i,0){\v};\fi\else\ifnum\i=5\ifnum\a=0\node[sort done]at(\i,0){\v};\else\node[sort cell]at(\i,0){\v};\fi\else\node[sort cell]at(\i,0){\v};\fi\fi\fi}
\end{tikzpicture}
\vspace{4mm}\only<1|handout:0>{29와 10: 역전이므로 swap}\only<2|handout:0>{29와 14: swap}\only<3|handout:0>{29와 37: 순서대로이므로 swap 없음}\only<4|handout:0>{37과 13: swap}\only<5|handout:1>{첫 pass 종료: 최댓값 37 확정}
\end{frame}
\begin{frame}{Bubble Sort의 시간과 성질}
\begin{tabular}{lccc}\toprule&best&average&worst\\\midrule 조기 종료 구현&$\Theta(n)$&$\Theta(n^2)$&$\Theta(n^2)$\\\bottomrule\end{tabular}
\vspace{4mm}\begin{itemize}\item stable: 오직 $A[i]>A[i+1]$일 때 swap하므로 같은 key는 지나치지 않는다.\item 조기 종료로 완전히 정렬된 입력은 $\Theta(n)$이다.\item nearly sorted 데이터의 실용적 적응성은 Insertion Sort와 다를 수 있다.\end{itemize}
\end{frame}
\begin{frame}{Bubble Sort Takeaway}
\sortingtakeaway{한 pass는 최대 원소 하나를 확정한다. 교환이 없으면 이미 정렬되었음을 증명한다.}
\begin{alertblock}{실무 관점}대부분의 범용 정렬에는 더 나은 선택이 있다. 하지만 invariant와 조기 종료를 학습하기 좋다.\end{alertblock}
\muted{앞의 두 알고리즘은 오른쪽 위치를 확정했다. Insertion은 반대로 정렬된 prefix를 한 칸씩 확장한다.}
\end{frame}
